\chapter{Every tool, and why}

Each tool gets a plain explanation, then what it does in this project, then why it was chosen over
the alternatives.

\section{COLMAP, used through pycolmap}

\begin{plainbox}
\textbf{What it is.} Software that works out the three dimensional shape of a scene from ordinary
two dimensional photographs, and where the camera was standing for each one. The technique is
called \emph{structure from motion}. It is the same idea your eyes use: two slightly different
views of the same thing let you judge depth.
\end{plainbox}

\begin{project}
It takes every frame from both videos and returns two things. First, the position and angle of the
camera for each frame. Second, a sparse cloud of three dimensional points marking corners and edges
it managed to recognise in several frames at once.
\end{project}

\textbf{Why this one.} It is the standard tool for the job and it is free. The Python package
\texttt{pycolmap} installs with pip, so there is no separate program to set up, which matters when
somebody else has to run your work. A robotics team will recognise it immediately.

\section{Depth Anything V2}

\begin{plainbox}
\textbf{What it is.} A neural network that looks at a single photograph and estimates how far away
every pixel is. It was trained on an enormous number of images, so it has learned the cues that
make things look near or far: perspective, relative size, and how objects overlap each other.
\end{plainbox}

\begin{project}
COLMAP returns a few hundred accurate points per frame, which is nowhere near enough to say whether
a chair is standing in a corner. This fills in a depth estimate for every single pixel.
\end{project}

\textbf{Why this one.} It runs on a laptop with no dedicated graphics card. The textbook
alternative, dense multi-view stereo, is more accurate but effectively requires a powerful GPU,
which would have made the project impossible to run or review on an ordinary machine.

\begin{watch}
Its output has no sense of scale. It knows the sofa is nearer than the wall; it does not know the
sofa is two metres away. Section~4.3 explains how that is fixed, and it is one of the more
interesting parts of the project.
\end{watch}

\section{CLIP}

\begin{plainbox}
\textbf{What it is.} A model trained on hundreds of millions of image and caption pairs. Give it a
picture and a list of candidate descriptions, and it scores how well each description fits. It was
never trained on a fixed list of objects, so the list can be changed without retraining anything.
\end{plainbox}

\begin{project}
Once the geometry has found a lump of space that changed, CLIP is shown a cropped photograph of
that space and picks the best match from a list of household objects, so the final report says
``a chair'' instead of ``cluster 3''.
\end{project}

\textbf{Why this one.} The geometry has already decided \emph{where} the object is, so a full
object detector is not needed, only a name. CLIP is small, runs on a laptop, and needs no training
data.

\section{DBSCAN, from scikit-learn}

\begin{plainbox}
\textbf{What it is.} An algorithm that finds groups in scattered data. Unlike most clustering
methods, you do not tell it how many groups to expect. It finds however many there are, and labels
loose scattered points as noise rather than forcing them into a group.
\end{plainbox}

\begin{project}
After the comparison there are thousands of individual cells of space marked as changed. DBSCAN
groups touching cells into objects, turning ``four thousand changed cells'' into ``a chair sized
lump here and a box sized lump over there''.
\end{project}

\textbf{Why this one.} Not knowing the number of objects in advance is the entire point, which
rules out anything that needs that number up front. Its willingness to discard scattered points as
noise is exactly what is wanted, because measurement error produces scattered points.

\section{OpenCV}

\begin{plainbox}
\textbf{What it is.} The standard library for working with images and video in code.
\end{plainbox}

\begin{project}
Reads the videos, pulls out frames, resizes them, and measures how blurry each one is so the bad
ones can be discarded. It also draws the synthetic test room described in Chapter~6.
\end{project}

\section{Plotly}

\begin{plainbox}
\textbf{What it is.} A charting library that produces interactive graphics which open in a web
browser.
\end{plainbox}

\begin{project}
Produces \texttt{scene.html}, the rotatable three dimensional model of the result.
\end{project}

\textbf{Why this one.} The output is a single file that opens in any browser. The usual
alternative for 3D work, Open3D, opens a desktop window instead, which is no use to somebody
reviewing the work from a link.

\section{Matplotlib}

\begin{plainbox}
\textbf{What it is.} The standard Python library for charts.
\end{plainbox}

\begin{project}
Draws the figures used in this document and in the project's README. They are generated from the
real output of a real run rather than drawn by hand, so they cannot quietly drift away from what
the code actually produces.
\end{project}

\section{NumPy}

\begin{plainbox}
\textbf{What it is.} The library that lets Python do arithmetic on millions of numbers quickly.
\end{plainbox}

\begin{project}
Everything numerical. The free space calculation described in Section~4.4 handles tens of millions
of points. In ordinary Python that would take hours. In NumPy it takes seconds.
\end{project}
